\section{Reading the fourth-order envelope equation}
\label{sec:dysthe}
\subsection{The equation in the paper's convention}
At $z=0$, Dysthe (2.19) is
\begin{align}
 2i\left(A_t+\frac12A_x\right)
 +\frac12A_{yy}-\frac14A_{xx}-A|A|^2
 &= -\frac{i}{8}(6A_{xyy}-A_{xxx})\notag\\
 &\quad+\frac{3i}{2}A(AA_x^*-A^*A_x)
       +\frac i2|A|^2A_x\notag\\
 &\quad+A(\Ph_x-i\Ph_z).
 \label{eq:dysthe-full}
\end{align}
Subscripts denote partial derivatives and $A_x^*=\partial_x(A^*)$.
The mean field obeys \eqref{eq:mean-full}, or its consistently reduced
version \eqref{eq:mean-bc}. The right side is fourth order under
Section~\ref{sec:asymptotics}'s ordering. The left side is the NLS balance.

\begin{center}
\begin{tabular}{@{}p{.34\linewidth}p{.59\linewidth}@{}}
\toprule Term & Role\\\midrule
$2i(A_t+A_x/2)$ & Envelope change following the group.\\
$A_{yy}/2-A_{xx}/4$ & Leading transverse and longitudinal dispersion.\\
$-A|A|^2$ & Leading nonlinear phase shift.\\
$-i(6A_{xyy}-A_{xxx})/8$ & Next Taylor terms of linear dispersion.\\
$\tfrac{3i}{2}A(AA_x^*-A^*A_x)+\tfrac i2|A|^2A_x$
& Nonlinear effects involving envelope gradients.\\
$A(\Ph_x-i\Ph_z)$ & Coupling to the induced mean response.\\
\bottomrule
\end{tabular}
\end{center}

The combination $AA_x^*-A^*A_x$ is imaginary. If $A=Re^{i\psi}$ with
$R>0$, it equals $-2iR^2\psi_x$. Its role therefore involves a local phase
gradient. The term $\partial_x|A|^2$ instead measures amplitude variation.
These structures need not have the same effect even though both are slow.

\subsection{An algebraic evolution form of the printed equation}
Divide \eqref{eq:dysthe-full} by $2i$, move the right side to the left,
and use $\Ph_z=(|A|^2)_x$. At the retained order this gives
\begin{align}
 A_t+\frac12A_x+\frac i8A_{xx}-\frac i4A_{yy}
 +\frac i2|A|^2A
 &-\frac1{16}A_{xxx}+\frac38A_{xyy}\notag\\
 &+|A|^2A_x-\frac14A^2A_x^*
 +\frac i2A\Ph_x=0.
 \label{eq:dysthe-evolution}
\end{align}
The coefficient change in the gradient terms uses the vertical mean-flow
boundary condition. For example,
$\tfrac12A\Ph_z=\tfrac12|A|^2A_x+\tfrac12A^2A_x^*$.
This is why comparing partially reduced equations can make coefficients
appear inconsistent. Always record which mean-field terms have been eliminated.

For a longitudinal profile, \eqref{eq:hilbert} turns the last term into
$-\tfrac i2A|D_x|(|A|^2)$ in this temporal potential-envelope convention.
Other variables and spatial evolution conventions can have different signs
and coefficients. The original paper and the later spatial free-surface
model are compared separately in Appendix~\ref{sec:project}.

\subsection{A source-consistency issue to recognize}
\label{sec:source-consistency}
There is an apparent coefficient inconsistency in the supplied scan.
Taken literally, the printed plus sign in front of
$i|A|^2A_x/2$ in (2.19) gives the evolution form
\eqref{eq:dysthe-evolution}. Linearizing it gives
\[
 P_1u+\frac34A_0^2u_x+P_2v=0,\qquad
 P_1v+\frac54A_0^2v_x-P_2u+A_0^2u=-\frac12\Ph_x.
\]
In contrast, the paper's printed (3.2), after using
$\Ph_z=2A_0^2u_x$, has $5A_0^2u_x/4$; its (3.3) has
$7A_0^2v_x/4$. Changing the displayed $+i|A|^2A_x/2$ to a minus
would produce those two coefficients. This is an algebraic consistency
observation about the supplied scan; no publisher erratum has been verified.
The literal intermediate (2.17) also does not reproduce all the gradient
coefficients of (2.19); see Appendix~\ref{sec:reduction}.

For learning the main stability argument, the next section starts explicitly
from the printed linearized system (3.2)--(3.3), which does reproduce
(3.6)--(3.13). We preserve (2.19) as printed rather than silently editing
it. You should not conclude that a failure to match these particular
coefficients is a gap in your prerequisite knowledge. Nor should this guide
be used as an unqualified numerical implementation specification.

\subsection{What the simplified system in the conclusion means}
Dysthe (4.1)--(4.3) keeps just NLS and the horizontal current:
\begin{equation}\label{eq:simplified}
 2i\left(A_t+\frac12A_x\right)
 +\frac12A_{yy}-\frac14A_{xx}
 =A(|A|^2+\Ph_x),
\end{equation}
with \eqref{eq:mean-bc}. It captures the leading improvement in the
\emph{growth-rate and instability-band} calculations emphasized in the
paper. It is not an equality between complete fourth-order dynamics.

In particular, the full equation also shifts the real perturbation
frequency by terms containing $\lambda A_0^2$ and third powers of
modulation wavenumber. Those terms matter for phase or transport even
when they do not change the growth rate at the retained order. Calling
the mean current dominant is a statement about the particular stability
observables being studied, not about every possible wave-packet calculation.

\paperlink{Read (2.19), p.~109, using the table above. Compare it with
the explicitly simplified system on p.~113.}
\checkpoint{Derive \eqref{eq:dysthe-evolution} from the printed
\eqref{eq:dysthe-full}. Then compare its linearization with the
source-consistency note, rather than assuming every printed line agrees.}
\clearpage
